%% ============================================================
%%  Chapter 1 — Executive Summary
%% ============================================================
\chapter{Executive Summary}
\label{ch:executive_summary}

\begin{warningbox}
  \fillme{One or two sentences stating what this report does and does not cover -- e.g. "this
  report identifies vulnerabilities only; no fixes were applied and no remediation changes were
  made to the codebase."}
\end{warningbox}

\vspace{0.5em}
\fillme{One paragraph summarising the overall posture: the categories of weakness found, and the
worst realistic outcome an unauthenticated or low-privilege attacker could achieve.}

%% ----- Severity tally -------------------------------------------------------
\section{Severity Tally}
\label{sec:severity_tally}

The assessment produced \textbf{\fillme{N} findings} (\fillme{F-001 through F-0NN}) across
\fillme{n} review passes.

\begin{table}[h!]
  \centering
  \caption{Finding count by severity}
  \label{tab:severity_tally}
  \begin{tabularx}{0.75\linewidth}{Xcc}
    \toprule
    \textbf{Severity} & \textbf{Count} & \textbf{Finding IDs} \\
    \midrule
    \rowcolor{sevCritical!15} Critical & \fillme{n} & \fillme{F-001, F-002, ...} \\
    \rowcolor{sevHigh!15} High & \fillme{n} & \fillme{F-003, ...} \\
    \rowcolor{sevMedium!15} Medium & \fillme{n} & \fillme{F-004, ...} \\
    \rowcolor{sevLow!15} Low / Info & \fillme{n} & \fillme{F-005, ...} \\
    \midrule
    \textbf{Total} & \textbf{\fillme{N}} & \\
    \bottomrule
  \end{tabularx}
\end{table}

%% ----- Top risks ------------------------------------------------------------
\section{Top Risks --- Fix First}
\label{sec:top_risks}

\fillme{One sentence introducing the list below.}

\begin{enumerate}[leftmargin=*, label=\textbf{\arabic*.}]
  \item \fillme{Headline risk cluster 1} (\Critical\ --- \S\ref{ch:example_category_one}, \fref{sec:F-001}).
  \item \fillme{Headline risk cluster 2} (\High\ --- \S\ref{ch:example_category_two}, \fref{sec:F-010}).
  \item \fillme{Headline risk cluster 3} (\Critical/\High\ --- \S\ref{ch:dependencies}).
\end{enumerate}

%% ----- Review structure -------------------------------------------------------
\section{Report Structure}
\label{sec:report_structure}

\begin{description}[leftmargin=2cm, style=nextline]
  \item[Part I (\S\S\ref{ch:example_category_one}--\ref{ch:dependencies})]
    Finding-by-finding write-up with descriptions, Proofs of Concept (illustrative), impacts,
    and mitigations, organized by vulnerability category.
  \item[Part II (\S\ref{ch:remediation})]
    Prioritized remediation roadmap (P0 immediate / P1 short-term / P2 hardening).
  \item[Appendix A] Complete finding matrix with ID, severity, confidence, and category.
  \item[Appendix B] Evidence file-and-class index mapping each category to its key source locations.
  \item[Appendix C] Runtime verification checklist for authorized testers to confirm/deny each finding.
\end{description}
